% PennyPilot rehearsal paper -- FROZEN SNAPSHOT, 14 October 2026.
% Do not edit: the working copy for the conference version is
% paper/rehearsal/main.tex. This snapshot reports completed experiments only.
% Numbers in Tables 1-2 are transcribed from paper/results/results.csv
% (scripts/paper/make_table.py over paper/results/raw/); numbers in
% Table 3 and Sec. 6 from paper/results/behavior.json
% (scripts/paper/behavior_table.py over paper/results/raw/*/behavior/).
%
% TODO(format): swap \documentclass/geometry for the target workshop's
% style file (e.g. \usepackage[preprint]{neurips_2025}) before submission.
\documentclass[10pt]{article}
\usepackage[margin=1in]{geometry}
\usepackage[T1]{fontenc}
\usepackage{amsmath,amssymb}
\usepackage{booktabs}
\usepackage{multirow}
\usepackage{xcolor}
\usepackage[numbers]{natbib}
\usepackage{hyperref}
\hypersetup{colorlinks=true,linkcolor=blue,citecolor=blue,urlcolor=blue}
\newcommand{\todo}[1]{\textcolor{red}{[TODO: #1]}}
\newcommand{\pmstd}[2]{#1\,{\scriptsize$\pm$#2}}

\title{Agents That Forget How to Chat:\\
Self-Distilled Rehearsal Prevents Tool-Call Collapse\\
Through SFT and RL Post-Training of a Shopping Agent}
\author{Param Madan\\Northeastern University\\\texttt{madan.pa@northeastern.edu}}
\date{}

\begin{document}
\maketitle

\begin{abstract}
Post-training a general instruction-tuned model into a tool-using agent can
silently destroy its ability to behave as a general assistant. We study this
in a multi-turn, permission-gated shopping environment with hidden user
constraints, fine-tuning Qwen2.5-7B-Instruct with LoRA, first by supervised
fine-tuning (SFT) on expert demonstrations and then by RLOO reinforcement
learning against a ground-truth reward. A specialist trained only on shopping
data reaches \pmstd{93.8}{2.7}\% task success on a held-out catalog after RL
(3 seeds), but answers only \pmstd{22.5}{9.0}\% of 40 general questions without
collapsing into store actions (base model: 95.0\%), and its IFEval
prompt-level strict accuracy falls from 64.5\% to \pmstd{38.8}{0.8}\%. MMLU
log-likelihood accuracy is unchanged (\(\approx\)74\% in every arm), and
generative probes point the same way: under the deployment prompt the specialist answers 0.0\% of
TriviaQA and multilingual math (MGSM) questions correctly on every seed, yet
with the system prompt removed the same post-RL checkpoints score
\pmstd{69.6}{1.2}\% and \pmstd{51.1}{0.8}\%, matching the base model.
This is strong evidence of prompt-conditioned behavioral interference rather
than measurable knowledge degradation. Mixing 25\%
\emph{self-distilled} rehearsal (open-ended Dolly prompts answered by the
untrained base model) into SFT keeps both measures at or above the base model through SFT and RL
(100.0\% Retention-40, \pmstd{65.5}{0.5}\% IFEval, \pmstd{69.9}{0.2}\%
TriviaQA under the deployment prompt). Task success is lower on
average (\pmstd{83.9}{9.4}\% vs.\ \pmstd{93.8}{2.7}\%) and more variable
across seeds, although with three seeds this gap is not statistically
significant (exact permutation test \(p=0.30\)). We release code, configurations, per-run manifests and
all raw evaluation outputs.
\end{abstract}

\section{Introduction}

A common recipe for building an LLM agent is to take an instruction-tuned
model, fine-tune it on task demonstrations, and then improve it with
reinforcement learning against a task reward. Deployed agents, however, rarely
see only in-task traffic: a shopping assistant is also asked what the capital
of Australia is. If post-training teaches the model to respond to
\emph{everything} with a store action, the agent becomes worse at the general
assistance it was supposed to retain.

Catastrophic forgetting under sequential training is long
known~\citep{mccloskey1989catastrophic,kirkpatrick2017overcoming} and has been
measured in LLM continual fine-tuning~\citep{luo2023empirical}; replaying old
data (rehearsal) is one of the oldest mitigations~\citep{robins1995catastrophic},
and mixing pre-training data into RLHF updates is standard
practice~\citep{ouyang2022training}. Closest to our method, self-distillation fine-tuning
(SDFT)~\citep{yang2024selfdistillation} reduces forgetting by having the seed
model rewrite the task data in its own words before SFT. Our work differs in
three respects. (i)~\emph{Setting}: we study a tool-using agent trained by SFT
\emph{followed by} RL, and ask whether the RL stage repairs, compounds or
leaves unchanged what SFT did. (ii)~\emph{Failure mode}: the loss we measure
is not degraded knowledge or safety but \emph{tool-call collapse}, where the
agent answers off-task questions with store actions; we show that a
knowledge benchmark misses it entirely. (iii)~\emph{Mechanism and cost}: the
task data is left untouched and the base model's answers to unrelated
open-ended prompts are mixed in as a separate rehearsal stream, and we report
what this costs in task success. This paper contributes a small, controlled,
fully reproducible measurement:

\begin{itemize}
  \item \textbf{Behavioral interference, not measurable knowledge loss.} Shopping-only
  post-training collapses off-task chat into store actions and halves IFEval
  accuracy, while MMLU (scored by log-likelihood, independent of output
  format) is unchanged (\S\ref{sec:results}). Removing the system prompt
  restores the specialist's TriviaQA and MGSM accuracy to base-model level,
  which indicates a prompt-conditioned behavior rather than measurable loss of
  knowledge on these benchmarks (\S\ref{sec:behavior}).
  \item \textbf{RL neither causes nor repairs it.} For both arms, retention
  and IFEval after RL are close to their post-SFT values (within seed
  variation); the damage is set during SFT.
  \item \textbf{Self-distilled rehearsal prevents it, at a likely cost.}
  25\% rehearsal keeps Retention-40 and IFEval at base-model level on all
  three seeds. Post-RL task success is 9.9 points lower on average and more
  variable across seeds; with three seeds this gap is not significant
  (exact permutation test \(p=0.30\); Welch \(p=0.20\)), so we report it as a likely cost rather than an established
  one.
  \item \textbf{A contamination-controlled evaluation.} We found that the
  environment's originally registered evaluation split shared 39 exact hidden
  needs and 52 answer sets with the training pools; all results here use a
  held-out catalog with need-level de-duplication (\S\ref{sec:eval}).
\end{itemize}

\section{Environment}
\label{sec:env}

\paragraph{Task.} Each episode is a multi-turn conversation (at most 12 turns)
between the agent and a simulated shopper in a synthetic product catalog of
300 items. The shopper has a hidden budget and 3--4 hidden hard requirements
that are revealed only in response to a relevant clarifying question. The agent
acts by emitting one JSON action per turn (\texttt{ask\_user},
\texttt{search}, \texttt{select\_product},
\texttt{request\_cart\_permission}, \texttt{add\_to\_cart}) after a short
natural-language sentence. User turns in our evaluation are Spanish--English
code-switched.

\paragraph{Permission gate.} \texttt{add\_to\_cart} for a product succeeds only
if that exact product was previously granted by the shopper. The gate is
enforced in the environment's state transition, so a violation cannot be
rewarded.

\paragraph{Reward and success.} The episode reward is
\(R = 0.4\,q + 0.4\,a + 0.2\,p - 1.0\,v + \sum_t g_t\), where \(q\) is the
value quality of the purchased item against catalog ground truth, \(a\)
indicates the shopper accepted the item, \(p\) that permission was requested,
\(v\) a permission violation, and \(g_t\) a per-turn information-gain bonus for
questions that eliminate inconsistent products. There is no learned reward
model. An episode is a \emph{success} if \(q>0\) and \(v=0\). We additionally
report regret in US dollars against the cheapest item satisfying all hidden
requirements.

\section{Method}
\label{sec:method}

\paragraph{Model and adapters.} All arms fine-tune
Qwen2.5-7B-Instruct~\citep{qwen2025qwen25} with a rank-16
LoRA~\citep{hu2022lora}; the base weights stay frozen and only the adapter is
trained.

\paragraph{SFT.} 540 expert demonstrations recorded programmatically in the
environment on the training catalog (catalog seed 0), in mixed languages and
including cases where the shopper denies permission or no product satisfies
the needs, one epoch, learning rate
\(10^{-5}\), batch size 2, maximum length 2048 tokens. The loss is computed on
agent tokens only. The demonstration seed is fixed across runs; the training
seed controls initialization and batch order.

\paragraph{Self-distilled rehearsal.} We select 900 Dolly-15k~\citep{conover2023dolly}
prompts from the open-QA, general-QA, brainstorming and creative-writing
categories that have no context passage, after filtering out shopping-related
words, and answer each with the \emph{untrained} base model under a plain
helpful-assistant system prompt. Answers that hit the generation limit without
an end-of-sequence token are dropped (57 of 900), leaving 843 exemplars. The
rehearsal arm adds the first 180 exemplars to the 540 shopping demonstrations,
so rehearsal is 25\% of the SFT mixture. Because targets come from the same
base model, rehearsal asks the adapter to \emph{preserve} existing behavior
rather than to learn new content.

\paragraph{RL.} Starting from each arm's SFT adapter, we run 50 steps of
RLOO~\citep{ahmadian2024back} with \(k=8\) rollouts per prompt from a pool of
64 training scenarios, learning rate \(10^{-6}\), and a per-token k3 KL
penalty~\citep{schulman2020kl} (\(\beta=0.04\)) to the SFT adapter. The
advantage of rollout \(i\) is \(R_i - \frac{1}{k-1}\sum_{j\neq i} R_j\). The
policy and the reference are two LoRA adapters on one shared frozen base,
which fits a 7B policy-plus-reference on one 32\,GB V100. RL uses shopping
data only for both arms; rehearsal enters only through SFT.

\paragraph{Arms.} \emph{Specialist}: SFT on shopping only, then RL.
\emph{Rehearsal25}: SFT on the 75/25 mixture, then the identical RL stage.
Each arm is run with training seeds 1, 2, 3 and evaluated after SFT and after
RL.

\section{Evaluation}
\label{sec:eval}

\paragraph{Contamination audit and held-out split.} The environment's
previously registered evaluation split (scenario seed 5001, same catalog as
training) shares 39 exact hidden needs and 52 exact answer sets with the SFT
and RL pools; we verified this with a unit test. All shopping results here
instead use 128 scenarios on a catalog generated with an unseen seed (1),
dropping any scenario whose hidden need also appears in training. Every run is
evaluated on the same 128 scenarios.

\paragraph{General-ability metrics.} We measure three things, chosen to
separate \emph{knowing} from \emph{behaving}:
\begin{itemize}
  \item \textbf{Retention-40}: 40 fixed general questions (27 English, 13
  Spanish; disjoint in template from the rehearsal prompts), asked under the
  deployment system prompt with greedy decoding. A response counts as retained
  if it is natural-language prose of at least 8 characters and does not parse as
  a store action. This measures collapse into the agent format, not answer
  correctness.
  \item \textbf{IFEval}~\citep{zhou2023ifeval} prompt-level strict accuracy on
  the first 200 prompts, with the chat template and the deployment system
  prompt (generative; behavior).
  \item \textbf{MMLU}~\citep{hendrycks2021mmlu}, 0-shot, log-likelihood
  scoring, 10 items per subject (570 items; knowledge, independent of output
  format).
\end{itemize}
Both IFEval and MMLU use lm-evaluation-harness 0.4.5~\citep{gao2023lmeval}.

\paragraph{Behavior and correctness probes.} To separate \emph{whether} the
model answers from \emph{whether it is right}, and to test what the collapse
depends on, we run a second greedy evaluation on every SFT and RL checkpoint
of both arms and on the base model (\nolinkurl{scripts/paper/eval_behavior.py};
every generated answer is saved).
\begin{itemize}
  \item \textbf{Off-task tool rate}: the fraction of 240 non-shopping prompts
  (the 40 Retention-40 questions plus 200 held-out Dolly prompts, asserted
  disjoint from the rehearsal set) answered with a parseable store action.
  It is measured under four system prompts: the deployment prompt; the
  deployment prompt plus an explicit instruction to answer non-shopping
  questions directly in prose (\emph{deploy+direct}); a generic
  helpful-assistant prompt; and no system prompt.
  \item \textbf{Answer correctness}: 300 TriviaQA questions (no-context
  setting; correct if any normalized gold alias appears as a whole-word span
  of the reply) and 300 MGSM grade-school math problems, 50 each in Spanish,
  French, German, Russian, Chinese and Japanese (correct if the last number in
  the reply equals the answer), under the deployment prompt and with no
  system prompt, at most 256 new tokens. A reply that is a store action is
  scored as wrong.
  \item \textbf{Adapter-scale sweep}: the LoRA update is multiplied by
  \(\alpha\in\{0,0.25,0.5,0.75,1\}\) (\(\alpha=0\) is the base model) and the
  tool rate is measured on the 40 Retention-40 questions under the deployment
  prompt.
\end{itemize}
We
report the mean and sample standard deviation over three seeds and, per run,
Wilson 95\% intervals on task success. Arm comparisons use a Welch
\(t\)-test and an exact permutation test on the three per-seed values
(\nolinkurl{scripts/paper/stats.py}).

\section{Results}
\label{sec:results}

\begin{table}[t]
\centering
\caption{Main results on the held-out catalog (128 scenarios). Mean
\(\pm\) sample std over 3 training seeds; the base model is a single
deterministic run. Permission violations were 0.0\% for every trained run.
All values are percentages except regret.}
\label{tab:main}
\small
\begin{tabular}{llcccccc}
\toprule
& & \multicolumn{2}{c}{Shopping} & \multicolumn{3}{c}{General ability} \\
\cmidrule(lr){3-4}\cmidrule(lr){5-7}
Arm & Stage & Success \(\uparrow\) & Regret \$ \(\downarrow\) & Retention-40 \(\uparrow\) & IFEval \(\uparrow\) & MMLU \\
\midrule
Base 7B & -- & 0.0 & -- & 95.0 & 64.5 & 74.0 \\
\midrule
\multirow{2}{*}{Specialist} & SFT & \pmstd{64.8}{7.5} & \pmstd{0.61}{1.06} & \pmstd{19.2}{7.2} & \pmstd{39.7}{0.8} & \pmstd{73.4}{0.1} \\
 & SFT\(\to\)RL & \textbf{\pmstd{93.8}{2.7}} & \pmstd{0.00}{0.00} & \pmstd{22.5}{9.0} & \pmstd{38.8}{0.8} & \pmstd{73.5}{0.2} \\
\midrule
\multirow{2}{*}{Rehearsal25} & SFT & \pmstd{51.6}{30.4} & \pmstd{0.63}{1.09} & \textbf{\pmstd{100.0}{0.0}} & \textbf{\pmstd{66.0}{0.5}} & \pmstd{74.0}{0.4} \\
 & SFT\(\to\)RL & \pmstd{83.9}{9.4} & \pmstd{0.03}{0.06} & \textbf{\pmstd{100.0}{0.0}} & \pmstd{65.5}{0.5} & \pmstd{74.0}{0.4} \\
\bottomrule
\end{tabular}
\end{table}

\paragraph{The base model cannot do the task.} Under the deployment system
prompt the untrained model never produces a valid store action on this split
(0\% success, 100\% invalid actions), so all task success in
Table~\ref{tab:main} comes from post-training.

\paragraph{Shopping-only post-training causes tool-call collapse.} The
specialist answers only 19.2\% of general questions as prose after SFT and
22.5\% after RL, versus 95.0\% for the base model. The failures are
qualitatively uniform: the model greets the user and emits a store action, for
example answering ``What's the capital of Australia?'' with a greeting followed
by \texttt{\{"action": "ask\_user", ...\}}. Retention is lower for Spanish
questions (15.4\% vs.\ 37.0\% English for specialist seed 1 after RL; single
seed). IFEval drops by roughly 25 points.

\paragraph{No measurable knowledge loss on MMLU.} MMLU accuracy is 73.4--74.0\% in every arm and
stage, against 74.0\% for the base model. Because MMLU is scored by
log-likelihood over answer options and is insensitive to output format, this
suggests that what changes is mainly the model's \emph{policy for when to act
as an assistant}; within the resolution of a 570-item subsample, we detect no
loss of the underlying knowledge. Practically, this means knowledge
benchmarks alone would have certified the specialist as unharmed.

\paragraph{RL does not change the measured general behavior.} For each arm, post-RL
retention changes by at most 3.3 points and IFEval by at most 0.9 points
relative to post-SFT, both comparable to the seed standard deviation.
In our setting (50 steps, KL-regularized to the SFT adapter, KL around
\(5\times10^{-4}\)--\(9\times10^{-4}\) per token in the logs we inspected), the
RL stage improves the task without moving Retention-40 or IFEval in either direction.

\paragraph{Rehearsal preserves Retention-40 and IFEval.} With 25\% rehearsal,
retention is 100.0\% on every seed and stage, and IFEval is
\pmstd{66.0}{0.5}\% after SFT and \pmstd{65.5}{0.5}\% after RL, slightly above
the base model's 64.5\%. Rehearsal targets are the base model's own answers,
so this is preservation rather than improvement; we do not interpret the
1-point difference.

\paragraph{The cost: lower and more variable task success.} After RL,
rehearsal25 reaches 83.9\% success against 93.8\% for the specialist.
Table~\ref{tab:seeds} shows the variance comes from individual seeds. After
SFT, rehearsal seed 3 succeeded on only 16.4\% of episodes, against 68.8--69.5\%
for seeds 1 and 2. That run never violated the permission gate, and every
purchase it made was optimal (zero regret); the remaining episodes ended without a
qualifying purchase (mean 11.8 turns vs.\ 11.1 for seed 1). RL
raised this seed to 74.2\%, but it remains the weakest run. Across seeds the post-RL success gap is not
statistically significant (Welch \(t\)-test \(p=0.20\); exact permutation
test \(p=0.30\)), whereas the gaps in Retention-40 (\(p=0.004\)) and IFEval
(\(p<10^{-5}\)) are (Welch). With three seeds per arm the exact permutation
test cannot go below \(p=0.10\); it reaches that floor for both general
metrics at both stages.

\begin{table}[t]
\centering
\caption{Per-seed shopping success (\%) with Wilson 95\% intervals
(\(n=128\) each).}
\label{tab:seeds}
\small
\begin{tabular}{llccc}
\toprule
Arm & Stage & Seed 1 & Seed 2 & Seed 3 \\
\midrule
Specialist & SFT & 68.8 [60.3, 76.1] & 69.5 [61.1, 76.8] & 56.3 [47.6, 64.5] \\
Specialist & RL  & 92.2 [86.2, 95.7] & 96.9 [92.2, 98.8] & 92.2 [86.2, 95.7] \\
Rehearsal25 & SFT & 69.5 [61.1, 76.8] & 68.8 [60.3, 76.1] & 16.4 [11.0, 23.8] \\
Rehearsal25 & RL  & 93.0 [87.2, 96.3] & 84.4 [77.1, 89.7] & 74.2 [66.0, 81.0] \\
\bottomrule
\end{tabular}
\end{table}

\section{What the Collapse Depends On}
\label{sec:behavior}

\begin{table}[t]
\centering
\caption{Behavior and correctness probes (\%). Off-task tool rate is over 240
non-shopping prompts; TriviaQA \(n=300\); MGSM \(n=300\) (6 languages). Mean
\(\pm\) sample std over 3 training seeds; the base model is a single
deterministic run.}
\label{tab:behavior}
\small
\setlength{\tabcolsep}{4pt}
\begin{tabular}{llccccccc}
\toprule
& & \multicolumn{3}{c}{Off-task tool rate \(\downarrow\)} & \multicolumn{2}{c}{TriviaQA \(\uparrow\)} & \multicolumn{2}{c}{MGSM \(\uparrow\)} \\
\cmidrule(lr){3-5}\cmidrule(lr){6-7}\cmidrule(lr){8-9}
Arm & Stage & deploy & deploy+direct & none & deploy & none & deploy & none \\
\midrule
Base 7B & -- & 9.2 & 5.4 & 0.0 & 35.7 & 68.3 & 7.0 & 48.3 \\
\midrule
\multirow{2}{*}{Specialist} & SFT & \pmstd{86.7}{7.8} & \pmstd{61.4}{8.8} & 0.0 & 0.0 & \pmstd{69.6}{1.4} & 0.0 & \pmstd{52.0}{0.9} \\
 & SFT\(\to\)RL & \pmstd{86.5}{10.3} & \pmstd{60.7}{6.9} & 0.0 & 0.0 & \pmstd{69.6}{1.2} & 0.0 & \pmstd{51.1}{0.8} \\
\midrule
\multirow{2}{*}{Rehearsal25} & SFT & \pmstd{0.1}{0.2} & \pmstd{0.1}{0.2} & 0.0 & \pmstd{70.0}{0.0} & \pmstd{70.8}{0.5} & \pmstd{53.8}{0.8} & \pmstd{52.4}{1.4} \\
 & SFT\(\to\)RL & \pmstd{0.1}{0.2} & \pmstd{0.1}{0.2} & 0.0 & \pmstd{69.9}{0.2} & \pmstd{70.7}{0.3} & \pmstd{53.9}{1.3} & \pmstd{53.4}{0.8} \\
\bottomrule
\end{tabular}
\end{table}

\paragraph{The specialist cannot answer under the deployment prompt.}
Under the deployment prompt the specialist answers 0.0\% of TriviaQA and MGSM
questions correctly on all six checkpoints (Table~\ref{tab:behavior}). Of
these 600 replies per checkpoint, \pmstd{88.2}{7.6}\% of TriviaQA replies and
\pmstd{86.0}{7.0}\% of MGSM replies are store actions after RL (88.1\% and
86.1\% after SFT), against at most 2.0\% on any rehearsal25 checkpoint. The
remaining specialist replies score zero as well,
and those we read redirect the user to shopping (``I'm here to help with PennyMart stuff! \ldots'')
rather than attempt an answer. On the 240 off-task prompts, 86.7\% (SFT) and
86.5\% (RL) of replies are store actions, against 0.1\% for rehearsal25
(Welch \(p=0.003\) after SFT, \(p=0.005\) after RL). With three seeds per arm
the exact permutation test cannot go below \(p=0.10\); it reaches that floor
for this metric and for TriviaQA and MGSM accuracy at both stages, since every
specialist seed is worse than every rehearsal25 seed. We do not report Welch
tests for the accuracy metrics, because the specialist scores 0.0\% on every
seed and the test is degenerate.

\paragraph{The answers are recoverable without the deployment prompt.} With the system prompt removed, the
same specialist checkpoints make no store actions and score
\pmstd{69.6}{1.2}\% on TriviaQA and \pmstd{51.1}{0.8}\% on MGSM after RL,
against 68.3\% and 48.3\% for the base model. Together with the unchanged MMLU
scores, this is strong evidence that post-training changed whether the
model answers under the deployment prompt rather than measurably degrading
what it knows; our benchmarks cover short factual and arithmetic answers
only, so subtler knowledge loss is not ruled out. Under a
generic helpful-assistant prompt the tool rate is also 0.0\% for every
checkpoint. An explicit instruction appended to the deployment prompt to
answer non-shopping questions directly lowers the specialist's off-task tool
rate only to about 61\%, so a prompt-level fix does not remove the behavior.

\paragraph{The collapse grows smoothly with the adapter.} Scaling the
specialist's LoRA update by \(\alpha=0, 0.25, 0.5, 0.75, 1\) gives tool rates
on the Retention-40 questions of 5.0, 14.2, 38.3, 63.3 and 82.5\% after SFT
and 5.0, 14.2, 38.3, 62.5 and 78.3\% after RL (means over three seeds; the
rate increases with \(\alpha\) on every seed). For rehearsal25 the rate is
5.0\% at \(\alpha=0\) (the base model) and 0.0\% at every \(\alpha>0\) on
every seed.

\paragraph{Rehearsal also preserves correctness.} After RL, rehearsal25 checkpoints
score \pmstd{69.9}{0.2}\% TriviaQA and \pmstd{53.9}{1.3}\% MGSM under the
deployment prompt, matching their no-prompt scores. We also
observe that the \emph{untrained} base model itself scores lower under the
deployment prompt (35.7\% TriviaQA, 7.0\% MGSM) than without it (68.3\% and
48.3\%). Reading its answers, we see it often decline off-topic questions as
outside the shopping role or ask a clarifying question instead of solving
the problem. The rehearsal checkpoints do not show this drop. A plausible
explanation is that the rehearsal targets were generated under a plain
assistant prompt, but this rests on a single deterministic base-model run
and we have not tested it.

\section{Status of the Results}
\label{sec:status}

This is an unpublished, not peer-reviewed manuscript describing an ongoing
study. It reports only completed experiments.
\begin{itemize}
  \item \textbf{Confirmed (3 training seeds, every run evaluated):} all
  results in Tables~\ref{tab:main}--\ref{tab:behavior}; the Retention-40 and
  IFEval differences between arms; the off-task tool-rate and correctness
  differences under the deployment prompt; that MMLU and no-prompt
  correctness show no measurable change after post-training; that RL leaves
  the measured general behavior unchanged.
  \item \textbf{Interpretation:} that the collapse is prompt-conditioned
  behavioral interference rather than knowledge loss. The measurements
  above strongly support it, but they cover MMLU, TriviaQA and MGSM only.
  \item \textbf{Preliminary:} the task-success cost of rehearsal (9.9 points;
  exact permutation \(p=0.30\), Welch \(p=0.20\); not significant with three
  seeds); the base model's lower correctness under
  the deployment prompt and its explanation (one run); the
  English--Spanish retention gap (one seed).
  \item \textbf{In progress and not reported here:} other rehearsal doses
  (10\% and 50\%); rehearsal with human-written instead of self-distilled
  answers; KL-to-base regularization during SFT as an alternative
  mitigation; training seeds 4 and 5; and replication on a second model
  family (Mistral-7B-Instruct-v0.3). No results from these runs appear in
  this version.
\end{itemize}

\section{Limitations}

\begin{itemize}
  \item \textbf{One model, one task, one rehearsal dose.} All results are for
  Qwen2.5-7B-Instruct with LoRA on a synthetic environment with a simulated
  shopper. We report a single rehearsal fraction (25\%); other doses and a
  second model family are in progress and not reported here (\S\ref{sec:status}).
  \item \textbf{Three seeds.} Seed variance in the rehearsal arm is large
  (std 30.4 points after SFT), so mean differences in task success should be
  read with caution.
  \item \textbf{Shared evaluation scenarios.} All runs use the same 128
  held-out scenarios; seeds vary training, not the test set. The Wilson
  intervals reflect scenario sampling for one run, not training variance.
  \item \textbf{Subsampled general benchmarks.} MMLU uses 570 items and IFEval
  200 prompts; small differences (\(\le\)1 point) are below resolution.
  \item \textbf{Retention-40 is a format metric.} It checks that the model
  answers in prose rather than with a store action; it does not grade answer
  correctness. TriviaQA and MGSM (\S\ref{sec:behavior}) cover correctness
  for short factual and arithmetic answers only; MGSM replies are capped at
  256 new tokens, which can truncate long solutions in every arm alike.
  \item \textbf{No comparison to other mitigations yet.} This version does
  not compare against KL-to-base during SFT or rehearsal with human-written
  answers (both in progress), or against weight-space regularization.
  \item \textbf{Short RL.} 50 RL steps with a small KL penalty; longer or more
  aggressive RL may interact with forgetting differently.
\end{itemize}

\section{Conclusion}

Turning a chat model into a tool-using agent through SFT and RL can make it
answer general questions with tool calls without measurably degrading its
knowledge, a failure that knowledge benchmarks do not detect. In our setting, the damage
is done during SFT, RL leaves it unchanged, and a 25\% mixture of the base
model's own answers to open-ended prompts prevents it on every seed as measured by Retention-40, IFEval
and answer correctness under the deployment prompt. Removing the system prompt
restores the specialist's correctness, which points to prompt-conditioned
behavioral interference rather than measurable knowledge degradation.
Task success was 9.9 points lower on average, a gap our three seeds cannot
distinguish from seed variance. Agent post-training pipelines should
report behavioral, not only knowledge, retention.

\paragraph{Reproducibility.} Code, the exact rehearsal prompt list, SLURM
scripts, per-run manifests (command line, git commit, host, library
versions) and all raw evaluation outputs are in the repository
(\nolinkurl{paper/results/raw/}); every number in this paper is generated from
them by \nolinkurl{scripts/paper/make_table.py} and
\nolinkurl{scripts/paper/behavior_table.py}. This snapshot is git tag
\texttt{snapshot-2026-10-14}.
Repository: \url{https://github.com/parammadan/pennypilot} (branch \texttt{paper/rehearsal}).

\bibliographystyle{plainnat}
\bibliography{refs}

\appendix
\section{Training on Volta GPUs}
\label{app:systems}
All runs use NVIDIA V100-SXM2-32GB GPUs, which lack native bf16 support and
FlashAttention. We train in fp16 autocast with dynamic loss scaling
(GradScaler) and fp32 trainable LoRA weights. Sharing one frozen base between
the policy and reference adapters avoids loading a second 7B model for the KL
term. RL steps took roughly 160--230\,s each in the logs we inspected; per
SLURM accounting, one rehearsal-arm SFT stage including its evaluation took
about 1.3\,h and one RL stage including its evaluation about 3.7\,h.
In total, the GPU jobs for this paper used 33.1 V100 GPU-hours by SLURM
accounting, including smoke tests and cancelled jobs.

\section{Additional baselines}
On the same held-out split, Qwen2.5-1.5B-Instruct reaches 2.3\% success
(5.5\% with shuffled search results), and Qwen2.5-7B-Instruct with an earlier,
more prescriptive system prompt reaches 1.6\%, confirming that the task is not
solvable zero-shot. General-ability metrics were measured only for the base 7B model under the deployment prompt (Table~\ref{tab:main}).

\end{document}
